\documentclass[11pt,a4paper]{amsart}
\usepackage[T1]{fontenc}
\usepackage{lmodern}
\usepackage{amssymb}
\usepackage{microtype}
\usepackage[hyphens]{url}
\usepackage[colorlinks=true,linkcolor=black,citecolor=black,urlcolor=black]{hyperref}
\usepackage[margin=2.5cm]{geometry}

\hypersetup{
  pdftitle={The ruler-and-compass plane is not 5-colourable},
  pdfauthor={Sergi Gal\'an},
  pdfsubject={The Hadwiger-Nelson problem},
  pdfkeywords={Hadwiger-Nelson problem, chromatic number of the plane, unit-distance graph, constructible numbers,
  formal verification}}

\newtheorem{theorem}{Theorem}
\newtheorem{corollary}[theorem]{Corollary}
\newtheorem{lemma}[theorem]{Lemma}
\theoremstyle{remark}
\newtheorem*{remark}{Remark}

\newcommand{\C}{\mathbb{C}}
\newcommand{\Q}{\mathbb{Q}}
\newcommand{\R}{\mathbb{R}}
\newcommand{\Z}{\mathbb{Z}}
\newcommand{\Ec}{E_{\mathrm{c}}}

\title{The ruler-and-compass plane is not 5-colourable}
\author{Sergi Gal\'an}
\thanks{These results were obtained with the assistance of an AI system (Claude, by Anthropic), which modified the
formal proof, ran the checks and helped to draft this note. The proof is OpenAI's \cite{OpenAI}; this note shows
that it works over a smaller field. No mathematician has checked the note yet. The author is not a professional
mathematician and submits it for independent verification.}
\date{Draft, 9 October 2026}
\subjclass[2020]{Primary 05C15, 52C10; Secondary 12F05, 37A15, 68V20}
\keywords{Hadwiger--Nelson problem, chromatic number of the plane, unit-distance graph, constructible numbers,
formal verification}


\begin{document}

\begin{abstract}
OpenAI recently proved that every colouring of the Euclidean plane with five colours gives the same colour to two
points at distance $1$, and formalized the proof in Lean. The proof works with the plane of algebraic points and
assumes that every algebraic number of modulus $1$ has roots of every order. We observe that square roots are
enough. Hence the plane of points constructible with ruler and compass is not $5$-colourable. These are the points
of the smallest subfield of $\C$ closed under square roots. By the de~Bruijn--Erd\H{o}s theorem, some finite
unit-distance graph with constructible vertices has no proper $5$-colouring. Two observations suffice: a
group-theoretic lemma about elements of order~$2$, and a choice of exponents in one multiple ergodic average that
needs no shift. We checked the modified proof in Lean~4. The kernel accepts it, and only the standard axioms are
used.
\end{abstract}

\maketitle

\section{The result}

The Hadwiger--Nelson problem asks for the least number $\chi(\R^2)$ of colours needed to colour the plane so that
points at distance $1$ get different colours. Isbell's hexagonal colouring gives $\chi(\R^2)\le7$, and de Grey
\cite{deGrey} proved $\chi(\R^2)\ge5$; see Soifer \cite{Soifer} for the history. OpenAI \cite{OpenAI} proved
$\chi(\R^2)\ge6$, with no regularity assumption on the colour classes. Their proof is not constructive, and we know
of no finite unit-distance graph without a proper $5$-colouring.

Identify the plane with $\C$. For a subset $S\subset\C$, a proper colouring of $S$ is a colouring in which points of
$S$ at distance $1$ get different colours. Let $\Ec$ be the field of \emph{constructible numbers}: the smallest
subfield of $\C$ closed under square roots, whose elements are the points constructible with ruler and compass from
$0$ and $1$ \cite[Ch.~10]{Cox}.

\begin{theorem}\label{thm:main}
$\Ec$ has no proper $5$-colouring.
\end{theorem}

So the unit-distance graph on $\Ec$ has chromatic number $6$ or $7$. Larger planes contain $\Ec$, so the theorem
also holds for them. Examples are the origami numbers (closed under square and cube roots \cite{Alperin}), the
numbers expressible by radicals, the algebraic numbers and $\C$ itself.

\begin{corollary}\label{cor:finite}
Some finite set of constructible points has no proper $5$-colouring.
\end{corollary}

\begin{proof}
A graph is $5$-colourable if all its finite subgraphs are, by the de~Bruijn--Erd\H{o}s theorem
\cite{deBruijnErdos}.
\end{proof}

The coordinates of such a set lie in a field obtained from $\Q$ by finitely many square roots. We know of no such
set. By contrast, the planes over $\Q(\sqrt2,\sqrt3)$ and $\Q(\sqrt3,\sqrt{11})$ have chromatic number $4$
\cite{P4C,Fischer}.

OpenAI's proof has two steps:
\begin{itemize}
\item \cite[Prop.~4.2]{OpenAI}, the forward half of their transfer theorem, turns a proper colouring of the plane into
a weak measurable colouring. Its proof first restricts the colouring to the countable plane $E$ of algebraic
points, and uses nothing else about it.
\item \cite[Thm.~1.4]{OpenAI} shows that the plane has no weak measurable $5$-colouring. It does not involve $E$.
\end{itemize}
We run the first step with $\Ec$ in place of the algebraic numbers. Section~\ref{sec:field} lists what that step uses
of its field. Sections~\ref{sec:involution} and~\ref{sec:triple} show that $\Ec$ has what it needs. Section~\ref{sec:lean}
describes the formal proof.

\begin{proof}[Proof of Theorem~\ref{thm:main}]
Suppose $\Ec$ had a proper $5$-colouring. By Sections~\ref{sec:field}--\ref{sec:triple}, the proof of
\cite[Prop.~4.2]{OpenAI} goes through with $\Ec$ in place of the algebraic numbers. It gives a weak measurable
$5$-colouring of the plane, and \cite[Thm.~1.4]{OpenAI} excludes one.
\end{proof}

\section{What the proof uses of the field}\label{sec:field}

OpenAI work with $E=\overline{\Q}\subset\C$, its real points $F=E\cap\R$, and the group $K=\{u\in E:|u|=1\}$
\cite[Def.~2.1]{OpenAI}. Their proof uses three properties of $E$:
\begin{enumerate}
\item[(E1)] $E$ is a countable subfield of $\C$, algebraic over $\Q$, closed under complex conjugation, and contains
$i$. Since $\Q(i)\subset E$, $E$ is dense in $\C$, $F$ in $\R$ and $K$ in the circle. These density facts are used
in \cite[Lemma~2.5, Lemma~4.1, Prop.~4.2]{OpenAI}.
\item[(E2)] $F$ is closed under square roots of its nonnegative elements. This is used in four places:
  \begin{itemize}
  \item the decomposition $x=k+\bar k$ with $k=x/2+i\sqrt{1-x^2/4}$ in \cite[Lemma~2.7]{OpenAI};
  \item the number $L=\sqrt{AB}$ in \cite[Lemma~3.3]{OpenAI};
  \item the points $t+i\sqrt{1-t^2}$ in \cite[\S3.5]{OpenAI};
  \item the moduli $|z|\in F$ of the elements $z\in E$.
  \end{itemize}
\item[(E3)] $K$ is divisible. It is used in three places:
  \begin{itemize}
  \item \cite[Lemma~2.4]{OpenAI} uses that a finite image of $K$ is trivial;
  \item \cite[Lemma~3.2]{OpenAI} uses that the power maps $u\mapsto u^n$ in its multiple averages are onto;
  \item \cite[Lemma~3.3]{OpenAI} uses directly the square map, in (3.15), and $u\mapsto u^{-2}$, in (3.13).
  \end{itemize}
\end{enumerate}
For $\Ec$, (E1) and (E2) hold, but (E3) holds only for the exponents $\pm2^k$.

The map $u\mapsto u^2$ is onto $K$: the square roots of $a+bi$, with $a^2+b^2=1$, have coordinates
$\pm\sqrt{(1+a)/2}$ and $\pm\sqrt{(1-a)/2}$. The map $u\mapsto u^3$ is not onto. Take
$u=(3+4i)/5=(2+i)/(2-i)$. It has valuation $1$ at the prime $2+i$ of $\Z[i]$, so it is not a cube in $\Q(i)$.
Hence $z^3-u$ is irreducible over $\Q(i)$, and its roots, of degree $3$ over $\Q(i)$, are not constructible.

\section{Elements of order two}\label{sec:involution}

In \cite[Lemma~2.4]{OpenAI}, a coset $d+C$ of positive mass is fixed by some rotation $u\ne1$. Then $(u-1)d\in C$,
and so $d\in C$. The formalization takes $u=-1$. That $-1$ fixes the coset follows from the next lemma, applied to
the action of $K$ on the finite orbit of the coset.

\begin{lemma}\label{lem:involution}
Let $G$ be a commutative group in which every element is a square, and let $f\colon G\to H$ be a homomorphism with
finite image. Then $f(g)=1$ for every $g\in G$ with $g^2=1$.
\end{lemma}

\begin{proof}
Squaring maps the finite group $f(G)$ onto itself, because every element of $G$ is a square. So squaring is
injective on $f(G)$. Since $f(g)^2=f(g^2)=1=1^2$, it follows that $f(g)=1$.
\end{proof}

The same lemma also gives the second use of divisibility in the formal proof. In the null-set statement for lines in
\cite[Lemma~2.4]{OpenAI}, it shows that $K$ is infinite: if $K$ were finite, the identity map of $K$ would send $-1$
to $1$.

\section{The triple average at the exponents \texorpdfstring{$-1,0,1$}{-1, 0, 1}}\label{sec:triple}

We use the notation of \cite[\S3]{OpenAI}:
\begin{itemize}
\item $X$ is an ergodic standard probability-preserving $K$-system, and $Y$ is a factor with $R=X\times_YX$ ergodic,
as in \cite[Lemma~3.2]{OpenAI};
\item $P=\mathbb{E}(\,\cdot\mid Y)$;
\item $U_u$ is the Koopman operator of $u\in K$. It is unitary, multiplicative, and commutes with $P$;
\item averages are taken along a F{\o}lner sequence of $K$.
\end{itemize}
The radial inequality $m(A)m(B)\le m(|A-B|)$ of \cite[Lemma~3.3]{OpenAI} rests on two facts:
\begin{itemize}
\item the identity $|A-Bu^2|=|-L+(A-B)u+Lu^2|$, for $L=\sqrt{AB}$ and $|u|=1$;
\item the multiple average of \cite[Lemma~3.2]{OpenAI} for the characters
$\psi_{-L}\cdot U_u\psi_{A-B}\cdot U_{u^2}\psi_L$, at the exponents $0,1,2$.
\end{itemize}
In the proof of \cite[Lemma~3.2]{OpenAI}, at (3.9), a common integer is added to the exponents of each term so that
none is zero. The formalization uses one shift for all terms and requires every shifted exponent to be onto. For
$0,1,2$, one of three consecutive nonzero integers is a nonzero multiple of~$3$, so this needs cube roots. A shift
chosen separately for each term would avoid them; we avoid the shift altogether.

Since $|u|=1$,
\[
|-L+(A-B)u+Lu^2|=|-Lu^{-1}+(A-B)+Lu|,
\]
and, since $U_u$ is unitary, multiplicative and commutes with $P$,
\[
\|P(a\cdot U_ub\cdot U_{u^2}c)\|=\|P(U_{u^{-1}}a\cdot b\cdot U_uc)\|.
\]
So it is enough to prove the conclusion~(3.2) of \cite[Lemma~3.2]{OpenAI} at the exponents $-1,0,1$, where the
middle factor is fixed.

\begin{lemma}\label{lem:triple}
Let $X$, $Y$ and $R$ be as above, and suppose that $u\mapsto u^2$ is onto $K$. For bounded functions $a,b,c$ on $X$,
\[
\operatorname{Avg}_u\bigl\|P(U_{u^{-1}}a\cdot b\cdot U_uc)-U_{u^{-1}}Pa\cdot Pb\cdot U_uPc\bigr\|_2^2\longrightarrow0
\]
along every F{\o}lner sequence.
\end{lemma}

\begin{proof}
Write each factor as $v=Pv+(v-Pv)$ and expand. The term with no centred factor is $U_{u^{-1}}Pa\cdot Pb\cdot U_uPc$.
It is $Y$-measurable, so $P$ fixes it.

Consider a term whose only centred factor is $b-Pb$. Its other two factors are $Y$-measurable, so its projection is
$U_{u^{-1}}Pa\cdot U_uPc\cdot P(b-Pb)=0$.

In every remaining term a centred factor sits at the exponent $-1$ or $1$. Let $W_u=U_{u^{-1}}a'\cdot b'\cdot U_uc'$
be such a term. By conditional independence,
\[
\|PW_u\|_2^2=\int_R\Phi\cdot G_u\,d\mu_R,\qquad \Phi=b'\otimes\bar b',\quad
G_u=U_{u^{-1}}(a'\otimes\bar a')\cdot U_u(c'\otimes\bar c').
\]
One of the two factors of $G_u$ is centred over $Y$, because $\mathbb{E}(w\otimes\bar w\mid Y)=|Pw|^2$.

Statement~(3.5) of \cite{OpenAI}, applied on $R$ to the two distinct nonzero exponents $-1$ and $1$, says that
$\operatorname{Avg}_uG_u\to0$ in $L^2(R)$. Its proof uses three facts:
\begin{itemize}
\item $R$ is ergodic;
\item the pair criterion~(3.3) holds on $R$ at $d=\pm1$;
\item the case $l=1$ holds through the power maps $u\mapsto u^{\pm2}$, which are onto.
\end{itemize}
The function $\Phi$ does not depend on $u$, so
$\bigl|\operatorname{Avg}_u\int_R\Phi\,G_u\,d\mu_R\bigr|\le\|\Phi\|_2\,\|\operatorname{Avg}_uG_u\|_2\to0$. This is
where it matters that the middle factor is fixed. There are finitely many terms, and the lemma follows.
\end{proof}

The power maps used are $u\mapsto u^{\pm1}$ and $u\mapsto u^{\pm2}$. The other two averages in
\cite[Lemma~3.3]{OpenAI} need only the square map:
\begin{itemize}
\item the pair at the exponents $0,2$ is the pair criterion~(3.3) with $d=2$;
\item the pair at the exponents $1,-1$ needs no shift, and its difference is $2$.
\end{itemize}
With Lemmas~\ref{lem:involution} and~\ref{lem:triple}, the rest of the proof of \cite[Prop.~4.2]{OpenAI} goes
through over $\Ec$ unchanged. This completes the proof of Theorem~\ref{thm:main}.

\section{The formal proof}\label{sec:lean}

OpenAI's formalization \cite{OpenAI} uses Lean~4 \cite{Lean4} with Mathlib \cite{Mathlib}; we used the same versions,
Lean~4.34.1 and Mathlib at commit \texttt{d13f23b7}. The proof is the import closure of their file
\texttt{Five.lean}: 70 modules of their library, on top of Mathlib. We changed it in two ways:
\begin{itemize}
\item a new file defines the field. It proves that the field is algebraic and closed under conjugation, contains $i$,
and satisfies~(E2). Countability and the square roots in $K$ are derived from these facts in other files. Eleven
files of the closure were adapted.
\item Two new lemmas in \texttt{Energy.lean} prove Lemma~\ref{lem:triple}, and the radial inequality uses it. Like
the rest of the formal proof, they are stated for invariant means rather than F{\o}lner sequences, so Lean proves the
invariant-mean form of Lemma~\ref{lem:triple}.
\end{itemize}
The final statement writes the field out, so the reader does not need to trust our definitions. In abbreviated form:
\begin{flushleft}\small
\texttt{theorem sqrtClosure\_plane\_not\_five\_colourable :}\\
\quad$\neg\,\exists\,c : (\mathtt{sInf}\ \{s : \mathtt{IntermediateField}\ \Q\ \C \mid
\forall x : \C,\ x^{2}\in s \to x\in s\}) \to \mathtt{Fin}\ 5,$\\
\qquad$\forall\,x\ y,\ \|(x:\C)-(y:\C)\| = 1 \to c\ x \ne c\ y$
\end{flushleft}
\verb|#print axioms| lists only \texttt{propext}, \texttt{Classical.choice} and \texttt{Quot.sound}. A second file
proves Corollary~\ref{cor:finite} from it, with Mathlib's compactness theorem for graph homomorphisms; the axioms are
the same.

We also checked the proof in two further ways:
\begin{itemize}
\item a fresh copy of OpenAI's sources (commit \texttt{fd4aeeb2}) was patched and built from scratch: the 72 modules
of the closure, namely their 70, the field file and the statement file, with none failing. Mathlib was prebuilt;
\item \texttt{leanchecker} replayed these 72 modules in the kernel. Mathlib was not replayed.
\end{itemize}
An AI-assisted review of an earlier version of these changes found no mathematical error.

Before the constructible case we built the origami and radical cases and checked their axioms. The origami case was
also rebuilt from scratch with the published script. The patches, the statements and a build script are in
\nolinkurl{lean/external/openai-five/fields/} of~\cite{Repo}, with a copy of OpenAI's licence (Apache~2.0).

\section{Remarks}

Property~(E2) says that $F$ is a Euclidean field, and the method needs it: with $B=1$,
\cite[Lemma~3.3]{OpenAI} uses $\sqrt A$ for every $A\ge0$ in $F$. Every Euclidean subfield of $\R$ contains the real
constructible numbers. So $\Ec$ is the smallest field of the form $F(i)$ to which the proof applies in this form.

Two questions remain open:
\begin{itemize}
\item Does the plane over the Pythagorean closure of $\Q$ have a proper $5$-colouring? There $K$ is still
$2$-divisible, because $(1\pm a)/2=((1\pm a)/2)^2+(b/2)^2$ is a sum of squares. So only (E2) fails.
\item How many square roots must be adjoined to $\Q$ to obtain a field whose plane has no proper $5$-colouring?
Corollary~\ref{cor:finite} says that finitely many suffice.
\end{itemize}
An explicit finite unit-distance graph with no proper $5$-colouring would bound the answer to the second question. It
would also give an explicit witness for $\chi(\R^2)\ge6$.

\begin{thebibliography}{99}

\bibitem{Alperin} R. C. Alperin, \emph{A mathematical theory of origami constructions and numbers}, New York J.
Math. \textbf{6} (2000), 119--133.
\bibitem{Cox} D. A. Cox, \emph{Galois Theory}, 2nd ed., Wiley, Hoboken, NJ, 2012.
\bibitem{deBruijnErdos} N. G. de Bruijn and P. Erd\H{o}s, \emph{A colour problem for infinite graphs and a problem
in the theory of relations}, Indag. Math. \textbf{13} (1951), 371--373.
\bibitem{deGrey} A. D. N. J. de Grey, \emph{The chromatic number of the plane is at least 5}, Geombinatorics
\textbf{28} (2018), 18--31; arXiv:1804.02385.
\bibitem{Fischer} K. G. Fischer, \emph{A planar geometric graph of chromatic number four}, Congr. Numer.
\textbf{104} (1994), 73--79.
\bibitem{Lean4} L. de Moura and S. Ullrich, \emph{The Lean 4 theorem prover and programming language}, Automated
Deduction -- CADE 28, Lecture Notes in Comput. Sci., vol. 12699, Springer, 2021, pp.~625--635.
\bibitem{Mathlib} The mathlib Community, \emph{The Lean mathematical library}, Proceedings of the 9th ACM SIGPLAN
International Conference on Certified Programs and Proofs (CPP 2020), ACM, 2020, pp.~367--381.
\bibitem{OpenAI} OpenAI, \emph{The Euclidean plane is not five-colorable}, preprint (23 September 2026), with a
formalization in Lean, \url{https://github.com/openai/math}, commit \texttt{fd4aeeb2}.
\bibitem{P4C} S. Gal\'an, \emph{A short proof that the planes over $\Q(\sqrt3,\sqrt{11})$ and $\Q(\sqrt2,\sqrt3)$
are $4$-chromatic}, version~6, 2026, not yet refereed, \texttt{papers/planes-4-chromatic/} in \cite{Repo}.
\bibitem{Repo} S. Gal\'an, \emph{The Hadwiger--Nelson problem over number fields}, code, data and formal proofs,
2026. Archived at Zenodo, \href{https://doi.org/10.5281/zenodo.22976635}{doi:10.5281/zenodo.22976635} (all
versions); source at \url{https://github.com/decalion89/chromatic-number-of-the-plane}.
\bibitem{Soifer} A. Soifer, \emph{The New Mathematical Coloring Book}, Springer, 2024.

\end{thebibliography}

\end{document}
